\chapter{Introduction}
\label{chap:introduction}

\section{Introduction to this Thesis}
The \emph{International Classification of Diseases (ICD)} is a globally recognized coding system maintained by the World Health Organization \citep{who2019icd11}, with its most recent iterations being ICD-10 and ICD-11. This system was developed to standardize the labeling of diseases, medical procedures, and underlying health conditions, thereby fostering interoperability across diverse healthcare environments. As a result, patient diagnoses and treatments are recorded more consistently, allowing healthcare providers to communicate more effectively, researchers to examine data with fewer ambiguities, and public health authorities to craft evidence-based policies.

\paragraph{}
When multiple ICD codes are assigned to a single patient’s record, the information becomes richer by including primary diagnoses and related comorbidities. Large-scale databases, such as the \emph{Medical Information Mart for Intensive Care-IV (MIMIC-IV)}, reflect this practice by covering not only the chief complaint but also secondary or rare conditions \citep{johnson2023mimicivscidata}. A major concern, however, arises if any of these codes are misassigned. Inaccuracies and omissions can ripple through subsequent clinical decision-making, research conclusions, and billing outcomes. Hence, the role of accurate coding becomes paramount in bridging the gap between raw clinical documentation and reliable medical insights.

\paragraph{}
From a financial standpoint, healthcare institutions depend on correct coding to claim reimbursements, and errors—either over-coding or under-coding—risk financial imbalances. Studies have highlighted the direct impact of clinical documentation on billing accuracy and institutional revenue, stressing the importance of precision in ICD assignments \citep{dong2022automated}. At a population level, aggregated ICD-coded data forms the bedrock of epidemiological tracking, policy development, and larger-scale medical research. Erroneous codes can alter the trajectory of these studies, potentially skewing public health interventions \citep{edin2023automated}.

\paragraph{}
Nevertheless, the manual process of ICD coding remains both tedious and vulnerable to human error. Coders routinely parse extensive and often unstructured discharge summaries that include technical jargon, acronyms, or shorthand specific to a given clinical context \citep{wrenn2010quantifying}. Rare or lesser-known conditions risk being misrepresented or overlooked in the daily rush of high-volume coding workloads. Automated or semi-automated tools have therefore been suggested as an avenue for alleviating these problems. Such solutions provide preliminary code assignments that human coders can validate, thus streamlining workflow and improving data quality \citep{campbell2020computer}.

\paragraph{}
Over the years, numerous computational methods have been devised to assist in the assignment of ICD codes. Early rule-based systems depended on dictionaries and handcrafted rules \citep{farkas2008automatic}, but these were limited in scalability and flexibility. Traditional machine learning techniques (e.g., support vector machines and logistic regression) made modest progress but often stumbled over the diverse, nuanced terminology of clinical text \citep{scheurwegs2017data,perotte2014diagnosis}. Momentum grew significantly with the advent of deep learning, starting with convolutional neural networks (CNNs) \citep{kim2014convolutional,mullenbach2018explainable} and recurrent neural networks (e.g., LSTMs \citep{hochreiter1997long}) that learned hierarchical or temporal patterns from unstructured notes. More recently, transformer-based models such as Bidirectional Encoder Representations from Transformers (BERT) have improved performance in many language tasks \citep{devlin2019bert}, though they often struggle with extremely long documents and heavily imbalanced data \citep{gao2021limitations,huang2022plm}.

\paragraph{}
Several specialized strategies have since been proposed to address these challenges, including extended context transformers, hierarchical embeddings, and curriculum learning over the ICD taxonomy \citep{nickel2017poincare,ren2022hicu,bengio2009curriculum}. These refinements aim to capture relevant details in lengthy clinical narratives and to handle the long-tail distribution in which a few codes dominate while a majority of rare codes receive minimal coverage \citep{rios2018few}. Alongside these methodological expansions, interpretability is receiving growing attention in automated ICD coding: clinicians must verify or scrutinize a model’s reasoning to ensure that each assigned code is justified by actual clinical evidence \citep{holzinger2017we,warner2024modernbert}. 

\paragraph{}
In parallel, new lines of research have begun focusing on *entity recognition* and *ontology-based approaches* to enhance the robustness of coding systems. Tools like \emph{cTAKES} \citep{savova2010ctakes} and \emph{MedCAT} \citep{kraljevic2021multi} identify medical concepts in text and map them to standardized vocabularies, such as SNOMED~CT \citep{donnelly2006snomedct} or the Unified Medical Language System (UMLS) \citep{bodenreider2004umls}. Although these pipelines do not automatically produce final ICD codes, they offer structured terminological insights that can reduce ambiguity in free-text data \citep{wu2018semehr}. More advanced systems incorporate ICD-specific knowledge directly into the classification process (e.g., hyperbolic embeddings of the ICD hierarchy \citep{cao2020hypercore} or label descriptions drawn from official manuals). This knowledge-aware approach has consistently shown promise in bridging the gap between unstructured narratives and the highly specialized code sets used in clinical practice.

\paragraph{}
An especially difficult hurdle is how to improve the recognition of rare or “zebra” codes. When diagnoses occur infrequently, data-driven models often lack sufficient positive examples for generalization \citep{rios2018few}. To combat this scarcity, researchers have begun to explore \emph{synthetic data generation}. By leveraging large language models (LLMs) to create artificial discharge summaries or augment existing text, the model’s exposure to underrepresented conditions can be increased \citep{falis2022horsestozebras,falis2024gpt35}. Ontologies like SNOMED~CT or the UMLS help ground these generated examples so that the synthetic narratives remain medically plausible \citep{kumichev2024medsyn}. For instance, a rare genetic disorder might have very few real cases in a training corpus; generating extra samples keyed to that condition’s known symptomatology can yield substantial performance gains, particularly in recall for low-frequency labels.

\paragraph{}
These generative methods build on a broader trend in natural language processing, in which powerful transformers (e.g., GPT-3.5 or GPT-4) are used to produce text that mimics human writing. Indeed, recent work suggests that LLMs prompted with curated medical knowledge can synthesize relatively realistic patient histories \citep{yang2023fewshoticd}. Although concerns persist about potential hallucinations or a lack of full clinical richness, experiments have shown that carefully designed prompts—possibly incorporating knowledge graphs or code descriptions—can yield synthetic notes beneficial for model training. Complementary approaches use synonyms or sibling diagnoses from an ontology to systematically “perturb” or augment existing text \citep{yuan2022codesynonyms}, thus expanding the model’s exposure to variant language for the same condition.

\paragraph{}
However, practical integration of these advanced methods requires consideration of real-world constraints. Data privacy regulations often limit the availability of large, richly annotated clinical corpora. Synthetic notes, if generated responsibly, can alleviate privacy concerns while still reflecting key features of real patient reports. Yet computational cost and resource constraints remain a challenge: training or fine-tuning long-context transformers can be expensive, and using LLM APIs extensively may be prohibitively costly for some healthcare institutions \citep{ji2021bert}. Techniques such as quantization, gradient checkpointing, or hierarchical batching \citep{geiping2023cramming} can reduce resource demands. Another paramount concern is reliability: automated coding systems must provide transparent justifications for the codes they produce if they are to earn the trust of clinicians. Interpretable mechanisms, such as label-specific attention or post-hoc attribution methods (e.g., Integrated Gradients, SHAP), facilitate manual review and verification \citep{teng2020explainable}.

\paragraph{}
To address these persistent challenges, this thesis proposes a framework that combines ontology-informed data augmentation, cutting-edge neural architectures, and interpretability techniques to enhance automated ICD coding. In particular, four major strategies anchor the work:

\begin{itemize}
    \item \textbf{Synthetic Data Generation for Rare Codes:} Large language models are leveraged in tandem with medical ontologies (e.g., SNOMED~CT, UMLS) to produce realistic training examples for less frequently occurring conditions \citep{kumichev2024medsyn,falis2022horsestozebras}. By enriching the dataset with synthetic discharge summaries, we strive to close the performance gap between common and rare codes.

    \item \textbf{Knowledge-Aware Multi-Label Architectures:} Neural models capable of handling thousands of code labels are examined, including transformer-based classifiers with extended input contexts. Techniques such as hyperbolic embeddings of the ICD hierarchy \citep{nickel2017poincare,cao2020hypercore} and label attention mechanisms \citep{vu2020ijcai} are integrated to improve performance and interpretability in multi-label classification.

    \item \textbf{Clinically Oriented Interpretability:} Interpretability remains crucial for adoption in clinical workflows. Methods like SHapley Additive exPlanations (SHAP), attention heatmaps, or concept-level attributions \citep{holzinger2017we,teng2020explainable} are incorporated to help coders and physicians quickly validate or question the system’s logic, especially for critical or rare diagnoses.

    \item \textbf{Resource-Conscious Training and Deployment:} Training strategies such as quantization and gradient checkpointing are used to mitigate heavy computational requirements \citep{geiping2023cramming}. This ensures that our proposed system remains accessible to healthcare organizations with limited GPU resources.
\end{itemize}

\paragraph{}
By uniting ontology-driven augmentation, state-of-the-art neural methods, and transparent validation, this work aims to significantly improve both the breadth and reliability of ICD coding systems. In so doing, we anticipate that more accurate, explainable coding will foster better clinical documentation, refine large-scale epidemiological analyses, and ultimately contribute to more effective patient care and public health interventions. The chapters that follow delve into the technical foundations, present empirical results, and discuss the real-world feasibility of deploying such a system in modern healthcare settings.
